\documentclass[12pt,a4paper]{article}
\usepackage[utf8]{inputenc}
\usepackage[vietnamese]{babel}
\usepackage[left=2cm,right=2cm,top=2cm,bottom=2cm]{geometry}
\usepackage{mathptmx}
\usepackage{enumitem}

\pagestyle{plain}
\linespread{1.15}

\begin{document}

% --- HEADER QUỐC HIỆU & CƠ QUAN BAN HÀNH ---
\noindent
\begin{minipage}[t]{0.45\textwidth}
    \begin{center}
        \small
        CÔNG AN THÀNH PHỐ HÀ NỘI\\
        \textbf{\underline{PHÒNG CẢNH SÁT HÌNH SỰ}}\\[0.1cm]
        Số: 11a/BB-KXNƠ\\
        \textit{V/v khám xét khẩn cấp chỗ ở đối tượng}
    \end{center}
\end{minipage}
\hfill
\begin{minipage}[t]{0.52\textwidth}
    \begin{center}
        \small
        \textbf{CỘNG HÒA XÃ HỘI CHỦ NGHĨA VIỆT NAM}\\
        \textbf{\underline{Độc lập - Tự do - Hạnh phúc}}\\[0.1cm]
        \textit{Hà Nội, ngày 25 tháng 07 năm 2026}
    \end{center}
\end{minipage}

\vspace{0.6cm}

% --- TIÊU ĐỀ VĂN BẢN ---
\begin{center}
    \textbf{\large BIÊN BẢN KHÁM XÉT KHẨN CẤP CHỖ Ở}\\
    \textit{(Phòng riêng của đối tượng Nguyễn Thanh Tùng tại khu tập thể Nhà máy Sứ)}
\end{center}

\vspace{0.3cm}

\noindent Căn cứ Lệnh khám xét khẩn cấp số 18/LKX-CSHS do Thủ trưởng Cơ quan CSĐT phê chuẩn;\\
Vào hồi 16 giờ 00 phút ngày 25 tháng 07 năm 2026, tại phòng trọ tầng 2, Khu tập thể Nhà máy Sứ, Phân khu Cảng, TP. Hà Nội.\\
Cán bộ chủ trì: Đại úy Lê Minh — Điều tra viên cùng sự chứng kiến của Tổ trưởng dân phố số 2 tiến hành khám xét chỗ ở của Nguyễn Thanh Tùng (Sinh năm 1990 — Nghề nghiệp: Công nhân xây dựng).

\section*{I. QUÁ TRÌNH KHÁM XÉT VÀ THU GIỮ ĐỒ VẬT}
Phòng trọ diện tích khoảng 15m2, đồ đạc đơn sơ. Trên nóc tủ quần áo bằng tôn cũ, cơ quan điều tra phát hiện và thu giữ **01 hộp các-tông niêm phong bụi bặm**, bên trong chứa các đồ vật sau:

\begin{enumerate}[label=\arabic*., leftmargin=1.5cm]
    \item \textbf{01 Cuốn sổ tay công nhân xây dựng:} Bìa da nâu, ghi chép lịch chấm công và ca trực công trường xây dựng (Không có dấu vết liên quan đến tội phạm — \textit{Vật chứng thông thường}).
    \item \textbf{Tập hóa đơn \& Biên lai cũ:} Gồm vài biên lai nộp tiền điện nước và hóa đơn sửa xe máy Wave BKS 29H1-829.12.
    \item \textbf{01 Chiếc đồng hồ đeo tay cũ:} Nhãn hiệu Orient cơ đã hỏng mặt kính và đứng kim.
    \item \textbf{Tập ảnh chụp lưu niệm:} 03 tấm ảnh chụp tập thể nhóm thợ xây công trình.
    \item \textbf{01 Bức tranh nhỏ đóng khung kính ố màu (VẬT CHỨNG THEN CHỐT):}
    \begin{itemize}[leftmargin=0.8cm]
        \item \textit{Mặt trước:} Nét vẽ màu sáp của trẻ con vẽ hình hai anh em đang nắm tay nhau mỉm cười dưới gốc cây.
        \item \textit{Mặt sau (Sau khi tháo nẹp khung):} Xuất hiện dòng chữ viết tay bằng bút mực tím đã phai màu theo năm tháng:
        \begin{center}
            \textit{"Thương em, anh xin lỗi vì đã không tìm thấy được em. An nghỉ em nhé. Hè 1998 — Anh Tùng"}
        \end{center}
    \end{itemize}
\end{enumerate}

\section*{II. ĐÁNH GIÁ Ý NGHĨA TRINH SÁT}
Dòng bút tích phía sau khung tranh xác nhận mối liên hệ ruột thịt sâu sắc giữa Nguyễn Thanh Tùng và bé Nguyễn Gia Huy (nạn nhân vụ ngạt khí tủ âm tường năm 1998). Dòng chữ thể hiện sự dằn vặt, ân hận tột cùng của Tùng vì đã không kịp thời tìm thấy em trai trong buổi chơi trốn tìm định mệnh năm 1998.

\vspace{0.8cm}

% --- CHỮ KÝ NGHỊ ĐỊNH 30 ---
\noindent
\begin{minipage}[t]{0.45\textwidth}
    \small
    \textbf{\textit{Nơi nhận:}}\\
    - Ban Chuyên án;\\
    - Viện KSMND TP. Hà Nội;\\
    - Lưu: VT, CSHS.
\end{minipage}
\hfill
\begin{minipage}[t]{0.5\textwidth}
    \begin{center}
        \small
        \textbf{CÁN BỘ CHỦ TRÌ KHÁM XÉT}\\
        \textit{(Ký, ghi rõ họ tên)}\\[1.5cm]
        \textbf{Đại úy Lê Minh}
    \end{center}
\end{minipage}

\end{document}
